\section{Related literature}\label{related-literature}

\subsection{2.1 Inflation-at-risk and distributional inflation
forecasting}\label{inflation-at-risk-and-distributional-inflation-forecasting}

Quantile regression provides a direct way to estimate conditional
distributional features without imposing a Gaussian predictive density
(\citeproc{ref-koenker1978quantiles}{Koenker and Bassett 1978}). The
asymmetric quantile or pinball loss is a strictly consistent scoring
function for a conditional quantile, so it supports coherent
out-of-sample comparison of upper-tail forecasts
(\citeproc{ref-gneiting2007proper}{Gneiting and Raftery 2007}).
Inflation-at-risk studies use this framework to quantify upside and
downside risks that are not visible in conditional-mean models.
López-Salido and Loria (\citeproc{ref-lopezsalido2024inflation}{2024})
show that the tails of the inflation distribution vary substantially
even when mean inflation appears stable. Banerjee et al.
(\citeproc{ref-banerjee2024inflation}{2024}) document nonlinear
exchange-rate and financial-condition effects in emerging economies.
Other work incorporates dynamic quantile regressions, time-varying
parameters, and Bayesian shrinkage (Busetti et al., 2015; Korobilis et
al., 2021).

This paper differs by focusing on a high-dimensional food-price panel, a
short three-month horizon, and real-time calibration failure after an
abrupt shift. The goal is predictive rather than causal. The forecast is
called Food Inflation-at-Risk because it is an upper conditional
quantile of future item-level food-price inflation; it is not a welfare
measure and should not be interpreted as the probability of a specific
macroeconomic event.

\subsection{2.2 Conformal calibration for dependent and shifting
data}\label{conformal-calibration-for-dependent-and-shifting-data}

Classical conformal prediction obtains finite-sample marginal coverage
under exchangeability, an assumption violated by dependent and
non-stationary time series. Adaptive conformal inference updates
calibration thresholds online to target long-run coverage under
distribution shift (\citeproc{ref-gibbs2021adaptive}{Gibbs and Candès
2021}). Later work develops local-regret adaptation, expert aggregation
over learning rates, and sequential residual-quantile estimation for
time series (Gibbs and Candès, 2024; Zaffran et al., 2022; Xu and Xie,
2023). Multi-step procedures explicitly recognise that a forecast error
may only become observable after the relevant horizon has passed
(\citeproc{ref-hallberg2024adaptive}{Hallberg Szabadváry 2024}).

The methods used here are conformal-style residual quantile calibrators
rather than exact conformal prediction sets. They operate on a single
upper-quantile forecast, use rolling or partially pooled residual
quantiles, and are evaluated by empirical coverage and pinball loss.
Dependence, source differences, and data-driven model development mean
that the forecasts do not carry a distribution-free conditional-coverage
guarantee. This distinction is retained throughout the paper.

\subsection{2.3 Expert aggregation and delayed
feedback}\label{expert-aggregation-and-delayed-feedback}

Prediction with expert advice studies how a learner can combine
forecasts and control regret relative to a reference expert. In
practical forecasting systems, feedback can be delayed because outcomes
are observed only after a reporting or forecast horizon. Joulani,
György, and Szepesvári (\citeproc{ref-joulani2013delayed}{2013}) show
that delay changes online-learning regret and requires explicit
information accounting. Expert aggregation has also been used to adapt
conformal procedures across learning rates and non-stationary
environments (\citeproc{ref-zaffran2022adaptive}{Zaffran et al. 2022}).

The controller in this paper is an empirical, domain-specific
implementation of this principle. It combines four transparent tail
forecasts. We do not claim a new universal regret bound. Instead, the
contribution is to enforce the exact three-month information delay,
estimate regime-conditioned expert priors before the structural break,
penalise unnecessary switching during validation, and evaluate dynamic
regret relative to an infeasible monthly hindsight oracle.

\subsection{2.4 Granular food-price data and
forecasting}\label{granular-food-price-data-and-forecasting}

Granular price systems create opportunities for disaggregated
forecasting and real-time monitoring. Recent evidence shows that
high-frequency and disaggregated price data can improve short-horizon
food-inflation nowcasts, although the value of additional data depends
on coverage and structural stability (\citeproc{ref-beer2026food}{Beer,
Ferstl, and Graf 2026}). Nigeria is particularly suitable for a granular
design because the National Bureau of Statistics publishes monthly
state-level prices for individual foods, fuel, and transport, while WFP
and the World Bank maintain market-level food-price systems. The NBS
Selected Food Price Watch is based on price collection across all states
and the FCT, while the World Bank Real-Time Food Prices system combines
direct observations with machine-learning estimates and is revised as
new information arrives (NBS, 2025; World Bank, 2026).

To preserve external validity, the main cross-source test uses
underlying WFP observations rather than treating World Bank imputed
values as independent ground truth. This choice sacrifices geographic
balance: the eligible post-2022 WFP sample is concentrated in Borno and
Yobe. The same-source NBS continuation and external WFP panel therefore
answer different questions and are interpreted jointly rather than
pooled as if they were interchangeable.
